\section{Making Change}\label{dp: sec: making change}\doHeader{Making Change}

\begin{mylist}
\item[\bf input:]
  A pair $I=(D, T)$ where
  $D=\{d_1, d_2, \ldots, d_k\}$
  is a set of \emph{denominations},
  and $T$ is a target
  % [review 2026-09-27] fix: denominations must be positive (d_i = 0 breaks the recurrence and the iteration); added "each d_i >= 1"
  (all non-negative integers, with each $d_i\ge 1$ and $1\in D$).

\item[\bf output:]
  The minimum size of any sequence $C=(c_1, \ldots, c_\ell)$
  s.t.~$c_i\in D$ for all $i$
  and $\sum_i c_i = T$.
\end{mylist}

\LN{greedy} described a greedy algorithm,
and proved that it worked for denominations $D=\{1,5,10,25\}$.
However, that algorithm doesn't work for all denominations.
For example, it fails for $D=\{1, 3, 4\}$ and target $T=6$.
Next we give a dynamic-programming algorithm that works in all cases.
Fix any input instance $(D=\{d_1, d_2, \ldots, d_k\}, T)$.

Our first goal is to come up with some divide-and-conquer
(recursive or inductive) understanding of the possible ways of making change for $T$.

To this end, partition the ways of making change for $T$ into up to $k$ types,
one for each denomination $d_i$ that could be used first,
which can be any denomination  $d_i$ that does not exceed $T$.
Among ways to make change for $T$ by first using a coin of denomination $d_i$,
the best uses that coin $d_i$, and is followed by the best possible
way of making change for $T-d_i$.
The best way of making change overall is the best among these (up to) $k$ types.
So, we could find the best way overall if we knew the best
way to make change for $T-d_i$, for each $d_i\le T$.

Applying this reasoning recursively,
for $t\in\{0,1,\ldots, T\}$, define $M(t)$ to be the minimum number of coins
needed to make change for $t$.
Then $1+M(t-d_i)$ is the minimum number of coins
needed to make change for $t$ if the first coin is $d_i$.
This reasoning gives the following recurrence relation:

\begin{lemma}\label{dp: lemma: change} $M(0) = 0$.  For $t\ge 1$,
  \(
    M(t) = \min\{ 1 + M(t-d_i) : d_i \in D,\, d_i \le t\}.
  \)
\end{lemma}
This recurrence relation naturally gives a recursive algorithm,
which we would memoize for efficiency.
Alternatively it gives the following iterative algorithm:

\begin{myframed}
  \begin{steps}
  \item[] \textsf{makeChange}$(D=\{d_1, d_2, \ldots, d_k\}, T)$:
    \step For $t=0,1,\ldots, T$:
    \step~~~$m[t] = \begin{cases}
      0 & \text{if } t=0 \\
      \min\{ 1 + m[t-d_i] : d_i \in D,\, d_i \le t\} & \text{otherwise.}
    \end{cases}$
    \step return $m[T]$
  \end{steps}
\end{myframed}

Correctness follows by the lemma and induction on $t$,
showing that $m[t]$ as computed by the algorithm equals $M(t)$
as defined earlier.
Each of the $T+1$ subproblems ($m[t]$ for each $t$)
takes time $O(k)$ to evaluate,
so the total time is $O(kT)$.

Let's summarize what we've figured out in a theorem:

\begin{theorem}
  There is an $O(kT)$-time algorithm for \prob{Making Change}.
\end{theorem}
\medskip

\pythonCode{dynamic_programming/making_change.py}

\subsection{Reducing to \prob{Shortest Path in a DAG}}\label{dp: sec: reducing to shortest path in a DAG}
Many dynamic-programming problems reduce to finding an optimal path in some implicit DAG.
For example, we can obtain the above algorithm for \prob{Making Change} this way.
Given an instance $(D, T)$ of  \prob{Making Change},
consider the following instance $(G=(V,E), s, t)$ of \prob{Shortest Path in a DAG}.

Take $V=\{0,1,\ldots, T\}$.  That is, there is a vertex for each subproblem $M(t)$.
For each $t$, add an edge from $t-d$ to $t$ for each $d\in D$ with $d\le t$.
So $E= \{(t-d, t) : t\in V; d\in D; d\le t\}$.
(There is an edge $(t-d, t)$
if subproblem $t-d$ occurs in the right-hand side of the recurrence
for subproblem $t$.)
Give each edge weight 1.


Then each path $p$ in $G$ from 0 to $T$
corresponds to a way of making change totaling $T$, and vice versa.
Specifically, if a path $p$ from 0 to $T$ goes through
a sequence of vertices such as
$0, c_1, c_1+c_2, \ldots, c_1+c_2+\cdots+c_n = T$,
then it corresponds to making change totaling $T$
using coins in $c=(c_1, c_2, \ldots, c_n)$.
The length of the path $p$ is the number of coins in $c$.
So, the minimum number of coins needed to make change totaling $T$
equals the distance $\dist 0 T$ from 0 to $T$ in $G$,
and another algorithm for \prob{Making Change} would work as follows:
\begin{myframed}
  \begin{steps}
  \item[]
    \textsf{makeChange2}$(D, T)$
    \comment{(return the min.~num.~of coins 
      needed to make change for $T$ with denominations in $D$)}
    
    \step construct the weighted DAG $G$ described above for $(D, T)$
    \step return \textsf{shortestPath}$(G, 0, T)$
  \end{steps}
\end{myframed}

This algorithm is essentially equivalent to the previous dynamic-programming algorithm
for \prob{Making Change}.

\pythonCode{dynamic_programming/dynamic_programming_paths.py}

How could you use this idea to solve the following problem?
\emph{Given $(D, T)$, compute the \emph{number of ways} to make change totaling $T$
using coins with denominations in $D$.}
(There is an ambiguity here: do we consider two ways of making change
different if they differ only in the order of their coins?
Is $c=(1, 1, 3)$ different than $c'=(3, 1, 1)$?
First answer the question assuming that order matters (so these $c$ and $c'$ are different).
Then solve it assuming order doesn't matter.
Hint: restrict to sequences $c$ such that, say $c_i \ge c_{i+1}$ for each $i$.)

% \paragraph{A long-form proof that the recurrence relation for {\sc Making Change}
%   is correct.}
%   Given an instance $(D,T)$ of \prob{Making Change},
%   recall that, for $t\in\{0,1,\ldots,T\}$,
%   $M[t]$ is defined to be the minimum number of coins needed
%   to make change totaling $t$, using denominations in $D$.
%   Recall the recurrence relation:
%   $M[0] = 0$, and, for $t\in\{1,\ldots,T\}$
%   \[
%     M[t] = \min\{ 1 + M[t-d] : d\in D, t-d\ge 0\}.
%   \]

% \begin{lemma}
%   The recurrence relation is correct.
% \end{lemma}
% \begin{longFormProof}
%   \step $M[0]=0$ because the empty sequence gives total 0.
%   \begin{block}[dp: t]
%     \step Consider any $t\in\{1,2,\ldots,T\}$.
%     \step 
%     First we show that $M[t]$ is \emph{at least} the right-hand side of the recurrence. 
%     \step
%     Let $c=(c_1, c_2, \ldots, c_n)$ be an optimal way of making change for $t$,
%     so $|c|=M[t]$.
%     \step Let $c' = (c_1, c_2, \ldots, c_{n-1})$ (the first $n-1$ coins in $c$).
%     \step There are $M[t]-1$ coins in $c'$, with total value $t-c_{n}$.
%     \step So it is possible to use $M[t]-1$ coins to make change with value $t-c_n$.
%     \step $M[t-c_n]$ is the minimum number of coins needed to make change with value $t-c_n$.
%     \step So $M[t]-1 \ge M[t-c_n]$.
%     \step That is, $M[t] \ge 1 + M[t-c_n]$. 
%     \step Since $c_n\in D$ and $t-c_n \ge 0$,
%     it follows that $M[t] \ge \min \{ 1 + M[t-d] : d\in D, t-d \ge 0\}$.
%     \step[dp: least] That is, $M[t]$ is at least the right-hand side of the recurrence. 

%     \smallskip
%     \step Next we show that $M[t]$ is \emph{at most} the right-hand side of the recurrence.
%     \begin{block}[dp: D]
%       \step Consider any $d\in D$ with $t-d\ge 0$.  We will show $M[t] \le 1 + M[t-d]$.
%       \step Let $c' = (c'_1, c'_2, \ldots, c'_m)$ be an optimal way of making change for $t-d$.
%       \step So $|c'| = M[t-d]$, and the total value of coins in $c'$ is $t-d$.
%       \step Let $c = (c'_1, c'_2, \ldots, c'_m, d)$ be obtained by adding a coin of denomination $d$ to $c'$.
%       \step There are $1+ M[t-d]$ coins in $c$, with total value $t-d + d = t$.
%       \step So it is possible to use $M[t-d]+1$ coins to make change with value $t$.
%       \step $M[t]$ is the minimum number of coins needed to make change with value $t$.
%       \step So $M[t] \le 1 + M[t-d]$.
%     \end{block}
%     \step By Block~\ref{dp: D}, for each $d\in D$ with $t-d\ge 0$,
%     we have $M[t] \le 1 + M[t-d]$.
%     \step That is, $M[t] \le \min\{1 +  M[t-d] : d\in D, t-d \ge 0\}$.
%     \step That is, $M[t]$ is at most the right-hand side of the recurrence.

%     \smallskip
%     \step By this and Step~\ref{dp: least}, $M[t]$ equals the right-hand of the recurrence.
%   \end{block}
%   \step By Block~\ref{dp: t}, for each $t\in\{1,\ldots,T\}$,  the recurrence holds.
% \end{longFormProof}

% \subsection*{External sources on Making Change}

% \begin{itemize}
% \item \JE Chapter 3 Exercise 1; DPV Exercises 6.17--6.19.
% \end{itemize}

%%% Local Variables:
%%% mode: latex
%%% TeX-master: "dynamic_programming"
%%% End:
